% Source: https://artofproblemsolving.com/community/c3h2415487_counting_in_a_circle
% Stable identity: aops-topic-2415487-q01
\textbf{(a)} Fix six positions around the circle. There are $6!=720$ ways to put the beads into the positions. The positions have $12$ symmetries: $6$ rotations (including turning by $0^\circ$) and $6$ reflections (flips). Applying the $12$ symmetries to one placement gives $12$ different placements: since the beads are all different, a symmetry that gives back the same placement must leave every position where it is, and only the identity does that. So each bracelet corresponds to exactly $12$ placements, and the number of bracelets is
\[\dfrac{720}{12}=\boxed{60}.\]
(Equivalently: fix one bead and arrange the other five in $5!=120$ ways around it, then divide by $2$ for flipping: $\dfrac{(6-1)!}{2}=60$.)

\textbf{(b)} Turning through $180^\circ$ and flipping over are different symmetries, so they give different placements. Read clockwise from the top:
\begin{itemize}
\item the start is $A,B,C,D,E,F$;
\item after turning through $180^\circ$ it is $D,E,F,A,B,C$;
\item after flipping over (about the vertical axis) it is $A,F,E,D,C,B$.
\end{itemize}

\begin{asy}
usepackage("amsmath");
size(9cm);
defaultpen(fontsize(10pt));
void hexa(pair c, string[] L, string cap){
 draw(circle(c,1),gray);
 for(int k=0;k<6;++k){pair p=c+dir(90-60k); dot(p); label(L[k],c+1.32*dir(90-60k));}
 label(cap,c+(0,-1.9));}
hexa((0,0),new string[]{"$A$","$B$","$C$","$D$","$E$","$F$"},"start");
hexa((4,0),new string[]{"$D$","$E$","$F$","$A$","$B$","$C$"},"turned $180^\circ$");
hexa((8,0),new string[]{"$A$","$F$","$E$","$D$","$C$","$B$"},"flipped over");
\end{asy}

A rotation keeps the clockwise order of the beads; a flip reverses it. (A $180^\circ$ turn of a flat figure is a reflection in a \emph{point}; flipping a bracelet over is a reflection in a \emph{line}, which is a different motion.) So all $12$ symmetries are different, and dividing by $6$ for the rotations and then by $2$ for flipping is correct.
